\chapter{量子力学初步}

\problem{2.1}{试求：
\begin{enumerate}
\item 电子、中子和氢原子在室温下的德布罗意波长；
\item 光子、电子、中子和氢原子具有 $0.1\ \mathrm{nm}$ 波长所对应的能量。
\end{enumerate}}
% layout
\answerarea

\problem{2.2}{试求：
\begin{enumerate}
\item 经过 $10\ \mathrm{kV}$ 电势差加速的电子束和质子束的德布罗意波长；
\item 在磁场强度为 $46\ \mathrm G$ 的均匀磁场中沿半径为 $0.5\ \mathrm{cm}$ 做圆周运动的电子的德布罗意波长；
\item 西欧中心的电子加速器 LEP 的电子能量为 $50\ \mathrm{GeV}$，求它的德布罗意波长。
\end{enumerate}}
% layout
\answerarea
\newpage

\problem{2.3}{证明在戴维孙和盖米的实验条件下不可能有与所测的一级极大值对应的二级和三级衍射峰。如果要得到二级衍射峰，并使它出现在 $50^\circ$ 处，需要用多大的加速电压？}
% layout
\answerarea

\problem{2.4}{在氯化钾晶体中，主平面间距为 $0.314\ \mathrm{nm}$，试比较能量为 $10\ \mathrm{keV}$ 的电子和光子在这些平面上的一级和二级布拉格反射方向与入射方向之间的夹角。}
% layout
\answerarea
\newpage

\problem{2.5}{中子衍射方法常用来进行结构分析，设有一窄束热中子射到晶体上，测得一级布拉格衍射角为 $30^\circ$，求晶体的布拉格面的间距。}
% layout
\answerarea

\problem{2.6}{用一束速度为 $400\ \mathrm{m\,s^{-1}}$ 的质子做杨氏双缝实验。如在缝后 $5\ \mathrm m$ 处的屏上观测到相距 $2\times10^{-4}\ \mathrm m$ 的明暗区，试求两缝的距离。}
% layout
\answerarea
\newpage

\problem{2.7}{试证明氢原子稳定轨道上正好能容纳下整数个电子的德布罗意波长，而且上述结果不但适用于圆轨道，也适用于椭圆轨道。}
% layout
\answerarea

\newpage

\problem{2.8}{一电子束在 $37\ \mathrm V$ 电场中加速并垂直通过一块有两狭缝 $A$ 和 $B$ 的板，缝宽 $0.1\ \mu\mathrm m$，两缝相距 $1\ \mu\mathrm m$，板后 $1\ \mathrm m$ 处有一个垂直于电子束方向的屏，屏上装有一位置灵敏探测器，能确定电子击中屏的位置。分别画出以下情况下到达屏上电子的相对位置分布图，并给以解释：
\begin{enumerate}
\item 缝 $A$ 开，缝 $B$ 关；
\item 缝 $B$ 开，缝 $A$ 关；
\item 缝 $A$、$B$ 都开；
\item 将斯特恩--盖拉赫装置连到板前，使只有 $s_z=\hbar/2$ 的电子能通过缝 $A$，$s_z=-\hbar/2$ 的电子能通过缝 $B$（参见下一章）；
\item 只有 $s_z=\hbar/2$ 的电子能通过缝 $A$，同时只有 $s_z=\hbar/2$ 的电子能通过缝 $B$；
\item 若束流强度减弱到如此之低，以至在某一时刻仅有一个电子通过装置，问有什么影响？
\end{enumerate}}
% layout
\answerarea
\newpage

\problem{2.9}{用不确定关系来解以下问题：
\begin{enumerate}
\item 设在非相对论情况下有 $H=p^2/(2\mu)-Ze^2/(4\pi\varepsilon_0r)$，求类氢离子的基态能量。
\item 设在相对论情况下有 $H=(p^2c^2+\mu^2c^4)^{1/2}-Ze^2/(4\pi\varepsilon_0r)$，求类氢离子的基态能量，此结果是否对所有 $Z$ 都正确？与（1）中结果有什么联系？
\item 一个质量为 $m$、动量为 $p$ 的粒子垂直入射到具有排斥势 $V_0$ 的表面的可穿深度。
\end{enumerate}}
% layout
\answerarea
\newpage

\problem{2.10}{试解释为什么在量子力学产生之前，一个很大的理论问题是如何防止原子发光，而在量子力学产生之后，一个很大的理论问题是如何使处于激发态的原子发光，什么使处于激发态的原子发光。}
% layout
\answerarea

\problem{2.11}{能量为 $1\ \mathrm{eV}$ 的电子束，其方向由两个相距 $1\ \mathrm m$ 的同样狭缝限定，为了最好地确定电子束方向，试问狭缝的宽度应为多少？}
% layout
\answerarea
\newpage

\problem{2.12}{一个电子被禁闭在一维盒子内，盒宽为 $10^{-10}\ \mathrm m$，电子处于基态，能量为 $38\ \mathrm{eV}$。计算：
\begin{enumerate}
\item 电子在第一激发态的能量；
\item 当电子处在基态时盒壁所受的平均力。
\end{enumerate}}
% layout
\answerarea

\problem{2.13}{\begin{enumerate}
\item 一个电子被禁闭在一个三维无限深势阱中，三个平行于 $x$、$y$ 和 $z$ 轴的边分别长为 $a$、$b$ 和 $c$，求电子可能有的波函数和能量。
\item 设这个势阱盒子是一立方体，边长为 $2\times10^{-10}\ \mathrm m$，求第一激发态相对基态的激发能量。
\end{enumerate}}
% layout
\answerarea
\newpage

\problem{2.14}{考虑质点在下列一维势中的运动：
\[
V(x)=\begin{cases}
\infty,&x<0,\\
0,&0\leq x\leq a,\\
V_0,&x>a.
\end{cases}
\]
证明束缚态能级由方程
\[
\tan\!\left(\frac{\sqrt{2mE}\,a}{\hbar}\right)=-\left[\frac{E}{V_0-E}\right]^{1/2}
\]
给出。不进一步求解，大致画出基态波函数的形状。}
% layout
\answerarea

\problem{2.15}{考虑势能为 $V(x)$ 的一维系统
\[
V(x)=\begin{cases}V_0,&x>0,\\0,&x<0,\end{cases}
\]
其中，$V_0>0$。若一能量为 $E$ 的粒子从 $x=-\infty$ 处入射，其透射率和反射率各为多少？考虑 $E$ 的所有能值。}
% layout
\answerarea
\newpage

\problem{2.16}{能量为 $1\ \mathrm{eV}$ 的电子入射到矩形势垒上，势垒高为 $2\ \mathrm{eV}$，为使穿透概率为 $10^{-3}$，求势垒的宽度。}
% layout
\answerarea

\problem{2.17}{在发现中子以前，人们曾经认为原子核由 $A$ 个质子和 $A-Z$ 个电子组成，试估算这种情况下的零点能，从而证明电子不可能被限定在像原子核这样小的区域内。}
% layout
\answerarea
\newpage

\problem{2.18}{试求氢原子基态的：
\begin{enumerate}
\item 波函数 $\psi$；
\item 坐标位置的平均值 $\bar{\boldsymbol r}$；
\item 电子沿径向分布的最大值处的 $r$ 值；
\item 势能平均值 $\bar V$；
\item 动能平均值 $\bar T$。
\end{enumerate}}
% layout
\answerarea

\newpage

